% Figure 2 -- pseudospectral landscapes.
% Drawn in PGFPlots from data/pseudo_*.dat (the stored fields) and
% data/pseudocont_*.dat (level sets extracted from those same fields).
% Every axis limit, colour limit and marker position below comes from
% data/pseudo_facts.tex, which scripts/tmlr_figdata.py writes next to the
% data, so the frame can never drift away from the field it encloses.
% The two panels share one window size and one square aspect ratio: their
% sharpness is meant to be compared, and that comparison is only honest if
% the plotting geometry is identical.
\begin{tikzpicture}

\pgfplotsset{
  pseudo/.style={
    width=6.45cm, height=6.45cm, scale only axis,
    view={0}{90},
    xlabel={$\Rey\omega$}, ylabel={$\Imy\omega$},
    label style={font=\small}, tick label style={font=\scriptsize},
    axis line style={draw=black!55, line width=0.3pt},
    tick style={black!55}, tick align=outside,
    colormap/viridis,
    colorbar horizontal,
    colorbar style={
      % far enough below the axis to clear the x label, which shares the gap
      height=0.19cm, at={(0,-0.205)}, anchor=north west, width=6.45cm,
      xtick style={draw=black!55},
      xticklabel style={font=\scriptsize},
      xlabel={$\log_{10}\sigma_{\min}/\sigma_{\max}$},
      xlabel style={font=\scriptsize, yshift=2pt},
      axis line style={draw=black!55, line width=0.3pt},
    },
  },
  % Level sets are drawn twice: a soft dark halo first, then the white line
  % on top.  A single white line disappears into the bright end of viridis
  % and a single dark line disappears into the dark end.
  contourhalo/.style={draw=black, line width=1.0pt, draw opacity=0.16,
                      no markers, unbounded coords=jump},
  contourover/.style={draw=white, line width=0.45pt, draw opacity=0.62,
                      no markers, unbounded coords=jump},
  trackedmode/.style={only marks, mark=*, mark size=1.9pt,
                      mark options={fill=white, draw=black, line width=0.6pt}},
  % the partner ring is drawn dark-then-light so it reads on either end of
  % the colour map, the same trick as the level sets
  partnerhalo/.style={only marks, mark=o, mark size=2.4pt,
                      mark options={draw=black, line width=1.5pt,
                                    draw opacity=0.35}},
  partnermode/.style={only marks, mark=o, mark size=2.2pt,
                      mark options={draw=white, line width=0.9pt}},
}
% node options resolve in /tikz/, so this one cannot live in \pgfplotsset
\tikzset{
  kappabox/.style={font=\scriptsize, text=white, inner xsep=2.6pt,
                   inner ysep=1.8pt, fill=black, fill opacity=0.55,
                   text opacity=1, rounded corners=1pt},
}

% ============================== (a) ordinary control ====================
\begin{axis}[pseudo,
  name=ctrl,
  xmin=\psAxmin, xmax=\psAxmax, ymin=\psAymin, ymax=\psAymax,
  xtick={1.5,1.7,1.9}, minor xtick={1.6,1.8,2.0},
  ytick={-0.6,-0.4,-0.2}, minor ytick={-0.5,-0.3,-0.1},
  point meta min=\psAzmin, point meta max=\psAzmax,
  colorbar style={xtick={-8,-7.5,-7,-6.5,-6,-5.5,-5,-4.5,-4,-3.5,-3,-2.5,-2}},
]
\addplot3[surf, shader=interp, draw=none,
          mesh/rows=\psGrid, mesh/cols=\psGrid]
  table[x=re, y=im, z=logsig] {data/pseudo_ordinary_control.dat};
\addplot[contourhalo] table[x=re, y=im]
  {data/pseudocont_ordinary_control.dat};
\addplot[contourover] table[x=re, y=im]
  {data/pseudocont_ordinary_control.dat};
\addplot[trackedmode] table[x=re, y=im]
  {data/pseudo_mark_ordinary_control.dat};
\node[kappabox, anchor=north west] at (rel axis cs:0.035,0.965)
  {$\kappa_T=18.3$};
\end{axis}

% ========================= (b) tightest interaction =====================
\begin{axis}[pseudo,
  name=tight, at={(ctrl.south east)}, xshift=1.95cm, anchor=south west,
  xmin=\psBxmin, xmax=\psBxmax, ymin=\psBymin, ymax=\psBymax,
  xtick={2.5,2.7,2.9}, minor xtick={2.6,2.8,3.0},
  ytick={-1.7,-1.5,-1.3}, minor ytick={-1.6,-1.4,-1.2},
  point meta min=\psBzmin, point meta max=\psBzmax,
  colorbar style={xtick={-8,-7.5,-7,-6.5,-6,-5.5,-5,-4.5,-4,-3.5,-3,-2.5,-2}},
]
\addplot3[surf, shader=interp, draw=none,
          mesh/rows=\psGrid, mesh/cols=\psGrid]
  table[x=re, y=im, z=logsig] {data/pseudo_tightest_interaction.dat};
\addplot[contourhalo] table[x=re, y=im]
  {data/pseudocont_tightest_interaction.dat};
\addplot[contourover] table[x=re, y=im]
  {data/pseudocont_tightest_interaction.dat};
\addplot[trackedmode] table[x=re, y=im]
  {data/pseudo_mark_tightest_interaction.dat};
\addplot[partnerhalo] table[x=re, y=im] {data/pseudo_pair.dat};
\addplot[partnermode] table[x=re, y=im] {data/pseudo_pair.dat};
\node[kappabox, anchor=north west] at (rel axis cs:0.035,0.965)
  {$\kappa_T=8.73\times10^{3}$};
\end{axis}

\node[anchor=south west, font=\bfseries\small, yshift=2pt]
  at (ctrl.north west)  {(a)\; ordinary control};
\node[anchor=south west, font=\bfseries\small, yshift=2pt]
  at (tight.north west) {(b)\; tightest stored interaction};
\end{tikzpicture}
